% Source base: Mathematics for Science Quiz, 2026 Science War Edition (2026-09-15)
% Selection principle: past-problem kernels + one standard step + high-yield famous results.

\finalchapter{1}{Set Theory and Logic}

\begin{formula}[Logical Equivalences]
바로 써야 할 논리 동치만 남긴다.
\[
(P\Rightarrow Q)\Longleftrightarrow(\neg Q\Rightarrow\neg P),
\qquad
\neg(P\Rightarrow Q)\Longleftrightarrow P\land\neg Q,
\]
\[
\neg(P\land Q)\Longleftrightarrow(\neg P)\lor(\neg Q),
\qquad
\neg(P\lor Q)\Longleftrightarrow(\neg P)\land(\neg Q),
\]
\[
\neg(\forall x\in A\, P(x))\Longleftrightarrow\exists x\in A\text{ s.t. }\neg P(x),
\quad
\neg(\exists x\in A\text{ s.t. }P(x))\Longleftrightarrow\forall x\in A\,\neg P(x).
\]
\end{formula}

\begin{theorem}[Cardinality Anchors]
\[
|\N|=|\Z|=|\Q|=\aleph_0,
\qquad
|(0,1)|=|\R|=\mathfrak c=2^{\aleph_0}>\aleph_0.
\]
유한집합 $A$에 대해 $|\Pow(A)|=2^{|A|}$. 따라서 $\Q$는 가산, $\R$는 비가산이다.
\end{theorem}

\begin{theorem}[Cantor and Schr\"oder--Bernstein]
Cantor 정리에 의해 모든 집합 $A$에 대해
\[
|A|<|\Pow(A)|.
\]
또한 $A\hookrightarrow B$와 $B\hookrightarrow A$인 단사함수가 모두 존재하면
\[
\boxed{|A|=|B|}
\]
이다 (Schr\"oder--Bernstein theorem).
\end{theorem}

\begin{fact}[Foundational Anchors]
\begin{itemize}
  \item \textbf{Russell's paradox}: 순진한 집합론의 모순을 보여 주는 대표 역설.
  \item \textbf{ZFC}: Zermelo--Fraenkel 집합론에 선택공리 (Axiom of Choice)를 더한 체계.
  \item 선택공리, Zorn's Lemma, Well-Ordering Theorem은 ZF 위에서 서로 동치.
  \item 연속체 가설 (CH): $\mathfrak c=\aleph_1$. G\"odel과 Cohen에 의해 ZFC로부터의 독립성이 확립되었다(일관성을 가정).
\end{itemize}
\end{fact}

\finalchapter{2}{Calculus and Analysis}

\finalsection{One-Variable Calculus}

\begin{formula}[Arithmetic, Geometric, and Power Sums]
\[
\sum_{k=1}^{n}k=\frac{n(n+1)}2,
\qquad
\sum_{k=1}^{n}k^2=\frac{n(n+1)(2n+1)}6,
\]
\[
\sum_{k=1}^{n}k^3=\left(\frac{n(n+1)}2\right)^2,
\qquad
\sum_{k=0}^{n-1}ar^k=a\frac{1-r^n}{1-r}.
\]
$|r|<1$이면
\[
\sum_{k=0}^{\infty}ar^k=\frac{a}{1-r}.
\]
\end{formula}

\begin{formula}[Asymptotic Anchors]
기본 극한보다 기억 가치가 큰 두 결과를 우선한다.
\[
\boxed{n!\sim\sqrt{2\pi n}\left(\frac ne\right)^n}
\qquad\text{(Stirling)}
\]
\[
H_n=1+\frac12+\cdots+\frac1n=\log n+\gamma+o(1).
\]
성장률 비교에서는 로그를 취하고, 계승이 보이면 Stirling 공식을 먼저 의심한다.
\end{formula}

\begin{formula}[Derivative Table: 실수하기 쉬운 항목]
너무 기본적인 항목보다 순간적으로 헷갈리기 쉬운 항목을 우선한다.
\[
(\arcsin x)'=\frac1{\sqrt{1-x^2}},\qquad
(\arccos x)'=-\frac1{\sqrt{1-x^2}},\qquad
(\arctan x)'=\frac1{1+x^2},
\]
\[
(\operatorname{arcsinh}x)'=\frac1{\sqrt{1+x^2}},\qquad
(\operatorname{arctanh}x)'=\frac1{1-x^2},
\]
\[
(\sinh x)'=\cosh x,\qquad
(\cosh x)'=\sinh x,\qquad
(\tanh x)'=\operatorname{sech}^2x.
\]
또한 $a^x$와 $\log_a x$에 대해
\[
(a^x)'=a^x\log a,\qquad
(\log_a x)'=\frac1{x\log a}.
\]
곱ㆍ몫ㆍ연쇄법칙은 별도 암기 항목이 아니라 즉시 쓰는 도구로 본다.
\end{formula}

\begin{formula}[Antiderivative Table: 암기 가치가 큰 형태]
\[
\int\frac{\dd x}{a^2+x^2}=\frac1a\arctan\frac xa+C,\qquad
\int\frac{\dd x}{\sqrt{a^2-x^2}}=\arcsin\frac xa+C,
\]
\[
\int\frac{\dd x}{x^2-a^2}=\frac1{2a}\log\left|\frac{x-a}{x+a}\right|+C,
\]
\[
\int\frac{\dd x}{\sqrt{x^2+a^2}}
=\operatorname{arcsinh}\frac xa+C
=\log\!\left|x+\sqrt{x^2+a^2}\right|+C',
\]
\[
\int \sec x\,\dd x=\log|\sec x+\tan x|+C,\qquad
\int \csc x\,\dd x=\log|\csc x-\cot x|+C.
\]
기본적인 거듭제곱ㆍ지수ㆍ삼각함수 적분은 자동화된 것으로 본다.
\end{formula}

\begin{formula}[Taylor and Maclaurin Series]
반드시 바로 떠올릴 다섯 개:
\[
e^x=\sum_{n=0}^{\infty}\frac{x^n}{n!},
\]
\[
\sin x=x-\frac{x^3}{3!}+\frac{x^5}{5!}-\cdots,
\qquad
\cos x=1-\frac{x^2}{2!}+\frac{x^4}{4!}-\cdots,
\]
\[
\log(1+x)=x-\frac{x^2}{2}+\frac{x^3}{3}-\cdots\quad(|x|<1),
\]
\[
\log(1-x)=-\sum_{n=1}^{\infty}\frac{x^n}{n}\quad(|x|<1),
\]
\[
(1+x)^\alpha=\sum_{n=0}^{\infty}\binom{\alpha}{n}x^n\quad(|x|<1).
\]
\end{formula}

\begin{theorem}[Fundamental Theorem of Calculus]
연속함수 $f$에 대해
\[
\frac{\dd}{\dd x}\int_a^x f(t)\,\dd t=f(x),
\qquad
\int_a^b f(x)\,\dd x=F(b)-F(a)\quad(F'=f).
\]
설명보다 이 두 형태를 바로 쓰는 것이 중요하다.
\end{theorem}

\begin{theorem}[Rolle and Mean Value Theorems]
적절한 연속성ㆍ미분가능성 가정 아래
\[
f(a)=f(b)\Longrightarrow\exists c\in(a,b):f'(c)=0,
\]
\[
\exists c\in(a,b):\quad f'(c)=\frac{f(b)-f(a)}{b-a}.
\]
조건과 결론만 기억한다.
\end{theorem}

\begin{formula}[Integration by Parts and Substitution]
\[
\int u\,\dd v=uv-\int v\,\dd u,
\qquad
\int_a^b f(g(t))g'(t)\,\dd t=\int_{g(a)}^{g(b)}f(x)\,\dd x.
\]
기법 자체보다 경계와 부호를 정확히 바꾸는 데 주의한다.
\end{formula}

\begin{formula}[Inverse-Function Derivative]
$y=f(x)$, $f'(x)\ne0$이면
\[
\boxed{(f^{-1})'(y)=\frac1{f'(x)}}.
\]
역함수를 직접 구하지 않고도 미분값을 계산할 수 있다.
\end{formula}

\begin{formula}[Gaussian Integral]
\[
\boxed{\displaystyle\int_{-\infty}^{\infty}e^{-x^2}\,\dd x=\sqrt\pi},
\qquad
\int_0^\infty e^{-ax^2}\,\dd x=\frac12\sqrt{\frac\pi a}\quad(a>0).
\]
정규분포, Gamma 함수, 극좌표와 연결되는 대표 결과이다.
\end{formula}


\begin{formula}[Beta--Gamma Toolkit]
\[
\Gamma(s)=\int_0^\infty x^{s-1}e^{-x}\,\dd x,
\quad
B(p,q)=\int_0^1x^{p-1}(1-x)^{q-1}\,\dd x,
\]
\[
\boxed{B(p,q)=\frac{\Gamma(p)\Gamma(q)}{\Gamma(p+q)}},
\qquad
\Gamma(n+1)=n!,\quad\Gamma\!\left(\frac12\right)=\sqrt\pi.
\]
또한
\[
\int_0^{\pi/2}\sin^m x\cos^n x\,\dd x
=\frac12B\!\left(\frac{m+1}{2},\frac{n+1}{2}\right).
\]
\end{formula}

\begin{formula}[Frullani's Integral]
적절한 가정 아래
\[
\int_0^\infty\frac{f(ax)-f(bx)}{x}\dd x
=(f(0)-f(\infty))\log\frac ba.
\]
형태가 보이면 먼저 $f(0)$와 $f(\infty)$를 확인한다.
\end{formula}

\finalheading{적분에서 먼저 확인할 것}
\begin{enumerate}
  \item 구간과 함수의 대칭: 홀함수/짝함수, $x\mapsto a-x$, $x\mapsto1/x$.
  \item 분자에 분모의 미분이 숨어 있는가.
  \item 치환 또는 부분분수가 즉시 되는가.
  \item $\sin^m x\cos^n x$, $x^{p-1}(1-x)^{q-1}$ 꼴이면 Beta/Gamma가 보이는가.
  \item 이상적분이면 수렴부터 확인한다.
\end{enumerate}

\begin{formula}[Series Tests]
\[
\sum_{n=1}^{\infty}\frac1{n^p}
\begin{cases}
\text{converges},&p>1,\\
\text{diverges},&p\le1.
\end{cases}
\]
비율판정법/근판정법에서 극한이 $<1$이면 절대수렴, $>1$이면 발산한다. 교대급수는 항의 크기가 $0$으로 감소하면 수렴한다.
\end{formula}

\begin{formula}[Two Useful Infinite-Series Anchors]
\[
\sum_{n=1}^\infty \frac{x^n}{n}=-\log(1-x)\quad(|x|<1),
\]
\[
\sum_{n=0}^\infty (-1)^n\frac{x^{2n+1}}{2n+1}=\arctan x\quad(|x|\le1\text{ 끝점에서는 조건부수렴}).
\]
특히 $x=1$에서 Gregory--Leibniz 급수
\[
\frac\pi4=1-\frac13+\frac15-\frac17+\cdots.
\]
\end{formula}

\begin{formula}[Wallis Product]
\[
\boxed{\displaystyle \frac\pi2=\prod_{n=1}^{\infty}\frac{(2n)^2}{(2n-1)(2n+1)}}.
\]
Wallis product는 $\pi$, Beta/Gamma, 삼각함수 적분과 연결되는 유명한 암기 결과이다.
\end{formula}


\begin{fact}[Basel Problem]
오일러 (Leonhard Euler)의 대표 결과:
\[
\boxed{\displaystyle\sum_{n=1}^{\infty}\frac1{n^2}=\frac{\pi^2}{6}}.
\]
과학퀴즈에서는 식 자체와 Euler를 함께 기억한다.
\end{fact}




\begin{fact}[Sophomore's Dream]
특이하고 외우기 좋은 적분-급수 항등식:
\[
\boxed{\displaystyle\int_0^1x^{-x}\,\dd x=\sum_{n=1}^{\infty}\frac1{n^n}}.
\]
\end{fact}

\begin{fact}[\(e\), \(\pi\), and an Open Trap]
에르미트 (Charles Hermite)는 $e$의 초월성을, 린데만 (Ferdinand von Lindemann)은 $\pi$의 초월성을 증명했다. 그러나
\[
e+\pi,\qquad e\pi
\]
각각이 무리수인지조차 알려져 있지 않다. 단, 둘이 동시에 대수적일 수는 없으므로 적어도 하나는 초월수이다.
\end{fact}

\begin{formula}[Fibonacci Numbers]
$F_0=0$, $F_1=1$, $F_{n+1}=F_n+F_{n-1}$이면
\[
F_n=\frac{\varphi^n-\psi^n}{\sqrt5},
\qquad
\varphi=\frac{1+\sqrt5}{2},\quad\psi=\frac{1-\sqrt5}{2},
\]
\[
\boxed{\sum_{n\ge0}F_nx^n=\frac{x}{1-x-x^2}},
\qquad
\frac{F_{n+1}}{F_n}\to\varphi.
\]
또한 $F_n$은 $\varphi^n/\sqrt5$에 가장 가까운 정수이다.
\end{formula}

\finalheading{수의 크기를 비교할 때}
거듭제곱 $a^b$와 $b^a$를 비교해야 하면 로그를 취해 $b\log a$와 $a\log b$를 비교한다. 특히
\[
f(x)=\frac{\log x}{x}
\]
는 $x=e$에서 최대이고 $x>e$에서는 감소한다. 2025 기출의 $11^9$와 $9^{11}$ 비교가 대표적이다.

\finalsection{Multivariable Calculus}

\begin{formula}[Gradient and Directional Derivative]
$f\colon\R^n\to\R$에 대해
\[
\nabla f=(\partial_1f,\ldots,\partial_nf)^{\mathsf T},
\qquad
D_{\mathbf u}f(\mathbf a)=\nabla f(\mathbf a)\cdot\mathbf u\quad(\|\mathbf u\|=1).
\]
Gradient는 가장 빠르게 증가하는 방향을 가리킨다.
\end{formula}

\begin{definition}[Jacobian and Total Derivative]
$F\colon\R^n\to\R^m$이면
\[
JF(\mathbf a)=\left(\frac{\partial F_i}{\partial x_j}(\mathbf a)\right).
\]
미분 가능하다는 것은
\[
F(\mathbf a+\mathbf h)=F(\mathbf a)+DF(\mathbf a)\mathbf h+\mathbf r(\mathbf h),
\qquad
\frac{\|\mathbf r(\mathbf h)\|}{\|\mathbf h\|}\to0
\]
인 최적 선형근사가 존재한다는 뜻이다.
\end{definition}

\begin{theorem}[Multivariable Chain Rule]
\[
D(G\circ F)(\mathbf a)=DG(F(\mathbf a))\,DF(\mathbf a).
\]
야코비 행렬의 곱 순서를 바꾸지 않는다.
\end{theorem}

\begin{formula}[Hessian and Second-Derivative Test]
\[
\Hess f=\left(\frac{\partial^2f}{\partial x_i\partial x_j}\right).
\]
임계점에서 Hessian이 양의 정부호이면 엄밀한 국소최소, 음의 정부호이면 엄밀한 국소최대, 부정부호이면 안장점이다. 판정 결과를 우선 기억한다.
\end{formula}

\begin{theorem}[Lagrange Multipliers]
제약식 $g(\mathbf x)=c$ 위의 정칙 극값에서는
\[
\boxed{\nabla f=\lambda\nabla g}.
\]
여러 제약식이면 각 gradient들의 선형결합을 쓴다.
\end{theorem}

\begin{formula}[Change of Variables]
좌표변환 $\mathbf x=F(\mathbf u)$에 대해
\[
\dd\mathbf x=\left|\det DF(\mathbf u)\right|\dd\mathbf u.
\]
극좌표/구면좌표에서는 Jacobian 인자를 빼먹지 않는다:
\[
\dd A=r\,\dd r\,\dd\theta,
\qquad
\dd V=\rho^2\sin\phi\,\dd\rho\,\dd\phi\,\dd\theta.
\]
\end{formula}


\begin{formula}[Vector Calculus Operators]
\[
\nabla f,
\qquad
\nabla\cdot\mathbf F,
\qquad
\nabla\times\mathbf F,
\qquad
\Delta f=\nabla\cdot\nabla f.
\]
$\mathbf F=\nabla f$이면 항상 $\nabla\times\mathbf F=\mathbf0$이다. 단일연결 영역에서는 적절한 정칙성 아래 curl이 0인 벡터장이 보존장이다.
\end{formula}

\begin{theorem}[Green, Stokes, Divergence]
형태를 알아보는 것이 핵심이다.
\[
\oint_{\partial D}P\dd x+Q\dd y
=\iint_D\left(\frac{\partial Q}{\partial x}-\frac{\partial P}{\partial y}\right)\dd A,
\]
\[
\int_{\partial S}\mathbf F\cdot\dd\mathbf r
=\iint_S(\nabla\times\mathbf F)\cdot\mathbf n\dd S,
\]
\[
\iint_{\partial V}\mathbf F\cdot\mathbf n\dd S
=\iiint_V\nabla\cdot\mathbf F\dd V.
\]
\end{theorem}

\finalsection{Real Analysis}



\begin{theorem}[Compactness in \(\R^n\)]
Heine--Borel:
\[
\boxed{K\subset\R^n\text{ compact}\iff K\text{ closed and bounded}.}
\]
연속함수는 콤팩트 집합에서 최대ㆍ최소를 갖고, 콤팩트 집합의 연속상은 콤팩트이다.
\end{theorem}

\begin{theorem}[Extreme and Intermediate Value Principles]
연속함수는 콤팩트 구간에서 최댓값과 최솟값을 가지며, 구간에서는 두 함수값 사이의 모든 값을 취한다. 계산보다 ``연속 + 콤팩트/구간''이라는 조건을 보는 것이 핵심이다.
\end{theorem}

\begin{formula}[Cauchy--Schwarz and Triangle Inequalities]
\[
|\langle \mathbf u,\mathbf v\rangle|\le\|\mathbf u\|\,\|\mathbf v\|,
\qquad
\|\mathbf u+\mathbf v\|\le\|\mathbf u\|+\|\mathbf v\|.
\]
Cauchy--Schwarz에서 등호가 성립하는 조건은 선형종속이다.
\end{formula}


\begin{fact}[Connectedness]
$\R$의 연결 부분집합은 정확히 구간이다. 연결집합의 연속상은 연결이고, 콤팩트 집합의 연속상은 콤팩트이다.
\end{fact}


\begin{fact}[Irrational Rotation]
Circle을 $\R/\Z$로 보고
\[
x\longmapsto x+\alpha\pmod1
\]
를 반복한다. $\alpha\in\Q$이면 모든 orbit은 주기적이고, $\alpha\notin\Q$이면 각 orbit은 원 위에 조밀하다. ``유리수 $\leftrightarrow$ 주기적, 무리수 $\leftrightarrow$ 조밀''를 한 쌍으로 기억한다.
\end{fact}

\begin{fact}[Cantor Set]
표준 Cantor set은 닫히고 유계이며 콤팩트ㆍ완전ㆍ내부가 조밀하지 않고 비가산이며 Lebesgue 측도 $0$이다. ``구간을 계속 제거했는데도 비가산''이라는 대비가 강한 quiz anchor이다.
\end{fact}

\begin{theorem}[Local Uniform Convergence: Useful Consequence]
함수열에서 미분과 극한의 순서를 바꾸려면 균등수렴에 가까운 조건을 확인해야 한다. 과거 기출처럼 ``국소 균등수렴 + 도함수 제어''이 주어지면 점별수렴만 보고 판단하지 않는다.
\end{theorem}

\finalsection{Complex Analysis}

\begin{formula}[Euler Formula and Roots of Unity]
\[
e^{i\theta}=\cos\theta+i\sin\theta,
\qquad
z^n=1\Longrightarrow z=e^{2\pi i k/n}\ (k=0,\ldots,n-1).
\]
주값 argument는
\[
\Arg z\in(-\pi,\pi]
\]
로 둔다.
\end{formula}

\begin{definition}[Holomorphic Function]
$f(z)=u(x,y)+iv(x,y)$가 복소미분 가능하면 정칙(holomorphic)이라 한다. 실전 판정에서 가장 먼저 보는 것은 Cauchy--Riemann 방정식이다.
\end{definition}

\begin{formula}[Cauchy--Riemann Equations]
$f=u+iv$가 충분히 매끄러울 때
\[
\boxed{u_x=v_y,\qquad u_y=-v_x}.
\]
복소미분 가능성의 대표 판정식으로 기억한다.
\end{formula}

\begin{theorem}[Cauchy's Integral Formula]
$f$가 적절한 영역에서 holomorphic이고 $a$가 단순 폐곡선 $\gamma$ 내부에 있으면
\[
f(a)=\frac{1}{2\pi i}\int_\gamma\frac{f(z)}{z-a}\dd z.
\]
더 일반적으로
\[
f^{(n)}(a)=\frac{n!}{2\pi i}\int_\gamma\frac{f(z)}{(z-a)^{n+1}}\dd z.
\]
\end{theorem}

\begin{theorem}[Residue Theorem]
\[
\int_\gamma f(z)\dd z
=2\pi i\sum \Res(f;a_k)
\]
($a_k$: $\gamma$ 내부의 고립특이점). 실수 이상적분을 contour integral로 바꾸는 문제의 핵심 결과이다.
\end{theorem}

\begin{theorem}[Rouch\'e's Theorem]
단순 폐곡선 $C$ 위에서
\[
|g(z)|<|f(z)|
\]
이면 $f$와 $f+g$는 $C$ 내부에서 중복도를 포함해 같은 개수의 영점을 갖는다. 복잡한 다항식의 unit-disk 영점 개수를 셀 때 강력하다.
\end{theorem}

\begin{theorem}[Liouville and Fundamental Theorem of Algebra]
유계 entire function은 상수함수이다. 그 결과를 이용하면 모든 비상수 복소다항식은 적어도 하나의 복소근을 갖는다.
\end{theorem}

\begin{formula}[Residues at Simple Poles]
$f(z)=g(z)/h(z)$이고 $h(a)=0$, $h'(a)\ne0$이면
\[
\Res(f;a)=\frac{g(a)}{h'(a)}.
\]
또한
\[
\Res\left(\frac{g(z)}{z-a};a\right)=g(a).
\]
\end{formula}


\begin{formula}[Principal Logarithm]
\[
\Log z=\log|z|+i\Arg z,
\qquad \Arg z\in(-\pi,\pi].
\]
일반 복소로그는 $\log z=\log|z|+i(\arg z+2\pi k)$처럼 다가치를 갖는다.
\end{formula}


\finalchapter{3}{Linear Algebra}

\begin{theorem}[Rank--Nullity]
$\mathbf A$가 $m\times n$ 행렬이면
\[
\rk\mathbf A+\dim\Null(\mathbf A)=n.
\]
열공간 차원과 행공간 차원, 즉 column rank와 row rank는 같다.
\end{theorem}

\begin{formula}[Subspace Dimension Formula]
부분공간 $U,V$에 대해
\[
\boxed{\dim(U+V)=\dim U+\dim V-\dim(U\cap V)}.
\]
특히 $U,V\subseteq\R^n$이면
\[
\dim(U\cap V)\ge\dim U+\dim V-n.
\]
\end{formula}

\begin{formula}[Determinant Essentials]
\begin{itemize}
  \item row swap: determinant의 부호가 바뀐다.
  \item 한 row에 $c$를 곱하면 determinant도 $c$배.
  \item 한 row에 다른 row의 multiple을 더해도 determinant는 변하지 않는다.
  \item 삼각행렬: $\det\mathbf A=\prod_i a_{ii}$.
  \item $\det(\mathbf A\mathbf B)=\det\mathbf A\det\mathbf B$.
  \item $\mathbf A$ invertible $\Longleftrightarrow \det\mathbf A\ne0$.
\end{itemize}
\end{formula}

\finalheading{Determinant를 계산할 때}
처음부터 cofactor expansion하지 않는다. 먼저 삼각/블록 구조, 0의 배치, 선형종속 또는 비례하는 행ㆍ열, 기본 행연산, 특수행렬 여부를 본다. Cofactor expansion은 0이 많을 때 특히 유리하다.


\begin{formula}[Determinant and Trace Identities]
\[
\det(\mathbf A^{\mathsf T})=\det\mathbf A,
\qquad
\det(\mathbf A^{-1})=\frac1{\det\mathbf A},
\qquad
\det(c\mathbf A)=c^n\det\mathbf A,
\]
\[
\tr(\mathbf A+\mathbf B)=\tr\mathbf A+\tr\mathbf B,
\qquad
\boxed{\tr(\mathbf A\mathbf B)=\tr(\mathbf B\mathbf A)}.
\]
더 일반적으로 cyclic permutation은 trace를 보존한다.
\end{formula}




\begin{fact}[LU Factorization and Pivoting]
가우스 소거법이 행 교환 없이 진행되면 $\mathbf A=\mathbf L\mathbf U$ 꼴이 자연스럽게 나온다. 하지만 모든 가역행렬이 피벗팅 없는 LU 분해을 갖는 것은 아니다. 일반적으로 순열행렬를 허용해
\[
\mathbf P\mathbf A=\mathbf L\mathbf U
\]
형태를 사용한다.
\end{fact}

\begin{theorem}[Equivalent Invertibility Conditions]
$n\times n$ 행렬 $\mathbf A$에 대해 다음은 동치이다.
\[
\det\mathbf A\ne0
\Longleftrightarrow
\rk\mathbf A=n
\Longleftrightarrow
\Null(\mathbf A)=\{\mathbf0\}
\Longleftrightarrow
0\notin\operatorname{spec}(\mathbf A).
\]
또한 정사각행렬에서
\[
\mathbf A\mathbf B=\mathbf I\Longrightarrow\mathbf B\mathbf A=\mathbf I.
\]
\end{theorem}

\begin{formula}[Characteristic Polynomial]
\[
p_{\mathbf A}(\lambda)=\det(\lambda\mathbf I-\mathbf A),
\qquad
p_{\mathbf A}(\lambda)=0\iff\lambda\text{는 고유값}.
\]
정의보다 trace, determinant, rank와 특수구조를 먼저 본다.
\end{formula}

\begin{formula}[Trace, Determinant, and Eigenvalues]
복소수까지 포함해 대수적 중복도로 세면
\[
\boxed{\sum_{i=1}^{n}\lambda_i=\tr\mathbf A},
\qquad
\boxed{\prod_{i=1}^{n}\lambda_i=\det\mathbf A}.
\]
$n\times n$ 행렬이 가질 수 있는 서로 다른 고유값의 최대 개수는 $n$이다.
\end{formula}

\finalheading{Eigenvalue 문제를 받으면}
\[
\boxed{
\tr\mathbf A
\;\to\;
\det\mathbf A
\;\to\;
\rk\mathbf A
\;\to\;
\text{obvious vector}
\;\to\;
\text{structure}
\;\to\;
p_{\mathbf A}(\lambda)
}
\]
특성다항식 계산은 마지막 수단에 가깝게 둔다.

\begin{formula}[Eigenvalues under Standard Operations]
$\mathbf A\mathbf v=\lambda\mathbf v$이면
\[
\mathbf A^k\mathbf v=\lambda^k\mathbf v,
\qquad
p(\mathbf A)\mathbf v=p(\lambda)\mathbf v.
\]
$\mathbf A$가 가역이면 $\mathbf A^{-1}$의 대응 고유값은 $\lambda^{-1}$. 또한 $\mathbf A+c\mathbf I$는 고유값을 $c$만큼 이동시킨다.
\end{formula}

\begin{theorem}[Cayley--Hamilton]
특성다항식 $p_{\mathbf A}(t)$에 대해
\[
\boxed{p_{\mathbf A}(\mathbf A)=\mathbf0}.
\]
고차 행렬 거듭제곱을 낮은 차수로 줄이거나 역행렬 식을 얻는 데 즉시 쓸 수 있다.
\end{theorem}




\begin{fact}[Nonzero Eigenvalues of \(AB\) and \(BA\)]
$\mathbf A$가 $m\times n$, $\mathbf B$가 $n\times m$이면 $\mathbf A\mathbf B$와 $\mathbf B\mathbf A$의 \textbf{0이 아닌 고유값은 대수적 중복도까지 같다}. 따라서 크기가 달라도 0이 아닌 스펙트럼 정보를 옮길 수 있다.
\end{fact}

\begin{fact}[Real Skew-Symmetric Matrices]
$\mathbf A^{\mathsf T}=-\mathbf A$이면 모든 실수 고유값은 $0$이고 0이 아닌 복소 고유값은 $\pm i\lambda$ 쌍으로 나타난다. 따라서
\[
\boxed{\rk\mathbf A\text{ is even}},
\qquad
n\text{ odd}\Longrightarrow\det\mathbf A=0.
\]
\end{fact}

\begin{formula}[Singular Value Decomposition]
모든 실수 $m\times n$ 행렬은
\[
\boxed{\mathbf A=\mathbf U\boldsymbol\Sigma\mathbf V^{\mathsf T}}
\]
로 쓸 수 있다. 특잇값(singular values)은 $\mathbf A^{\mathsf T}\mathbf A$의 고유값의 음이 아닌 제곱근이다.
\end{formula}

\begin{theorem}[Diagonalizability]
$\mathbf A$가 $n$개의 선형독립 고유벡터를 가지면
\[
\mathbf A=\mathbf P\mathbf D\mathbf P^{-1}
\]
으로 대각화 가능하다. 서로 다른 고유값에 대응하는 고유벡터는 선형독립이므로 고유값이 $n$개 모두 다르면 자동으로 대각화 가능하다.
\end{theorem}

\begin{theorem}[Spectral Theorem for Real Symmetric Matrices]
실수 대칭행렬은 직교대각화 가능하다.
\[
\mathbf A=\mathbf Q\mathbf D\mathbf Q^{\mathsf T},
\qquad
\mathbf Q^{\mathsf T}\mathbf Q=\mathbf I.
\]
모든 고유값은 실수이고, 서로 다른 고유값의 고유공간은 서로 직교한다.
\end{theorem}

\begin{formula}[Special Matrices]
\begin{center}
\small
\renewcommand{\arraystretch}{1.18}
\begin{tabularx}{0.96\textwidth}{@{}L{0.23\textwidth}Y Y@{}}
\toprule
구조 & 바로 쓸 결과 & 첫 반응 \\
\midrule
$\mathbf J_n$ & rank $1$; 고유값 $n,0,\ldots,0$ & $\ones$와 $\ones^\perp$ \\
$\mathbf J_n-\mathbf I_n$ & 고유값 $n-1,-1,\ldots,-1$ & 고유값 이동 \\
$\mathbf u\mathbf v^{\mathsf T}$ & rank $\le1$; 0이 아닌 고유값 $\mathbf v^{\mathsf T}\mathbf u$ & outer product \\
삼각행렬 & 대각성분이 고유값 & 특성다항식 계산 불필요 \\
직교행렬 & $\mathbf Q^{\mathsf T}\mathbf Q=\mathbf I$, $|\det\mathbf Q|=1$ & 길이ㆍ각도 보존 \\
사영행렬 $\mathbf P$ & $\mathbf P^2=\mathbf P$; 고유값 $0,1$ & 상/핵 분해 \\
멱영행렬 $\mathbf N$ & $\mathbf N^k=0$ for some $k$; 모든 고유값 $0$ & Jordan 연쇄 구조 \\
\bottomrule
\end{tabularx}
\end{center}
\end{formula}

\begin{formula}[Rank-One Matrix]
$\mathbf A=\mathbf u\mathbf v^{\mathsf T}$이면
\[
\mathbf A\mathbf u=(\mathbf v^{\mathsf T}\mathbf u)\mathbf u.
\]
따라서 $\mathbf v^{\mathsf T}\mathbf u\ne0$이면 이것이 유일한 0이 아닌 고유값이고
\[
\tr\mathbf A=\mathbf v^{\mathsf T}\mathbf u.
\]
\end{formula}


\begin{formula}[Sum of Outer Products]
\[
\mathbf A=\sum_{i=1}^{k}\mathbf u_i\mathbf v_i^{\mathsf T}
\quad\Longrightarrow\quad
\rk\mathbf A\le k.
\]
각 outer product의 rank가 최대 $1$이므로 rank의 subadditivity를 바로 적용한다.
\end{formula}

\begin{formula}[Projection]
0이 아닌 벡터 $\mathbf u$에 대한 직선 위 정사영:
\[
\operatorname{proj}_{\mathbf u}\mathbf v
=\frac{\mathbf v\cdot\mathbf u}{\mathbf u\cdot\mathbf u}\mathbf u.
\]
단위벡터 $\mathbf u$이면 정사영 행렬는
\[
\mathbf P=\mathbf u\mathbf u^{\mathsf T}.
\]
\end{formula}

\begin{fact}[Rotations in \(\R^3\)]
고유 회전행렬 $\mathbf R\in SO(3)$는 회전축 방향의 고유값 $1$을 가진다. 나머지 두 고유값은 일반적으로 $e^{\pm i\theta}$이다. 따라서
\[
\det\mathbf R=1.
\]
3D rotation 문제에서 ``실수 고유벡터 하나''는 회전축을 뜻한다.
\end{fact}

\begin{fact}[Pascal-Type Determinant]
이항계수로 이루어진 Pascal형 행렬은 행렬식이 $1$로 떨어지는 경우가 많다. 먼저 행ㆍ열 점화와 삼각분해를 찾는다.
\end{fact}

\begin{remark}
Jordan 표준형은 대각화할 수 없는 행렬에 남는 구조를 기술한다. 파이널에서는 세부 계산보다 ``대각화가 안 되어도 Jordan 표준형이 있다'' 정도만 기억한다.
\end{remark}


\finalchapter{4}{Differential Equations}

\begin{formula}[First-Order Exponential Growth and Decay]
\[
y'=ay
\quad\Longrightarrow\quad
y=Ce^{ax}.
\]
$a>0$이면 성장, $a<0$이면 감쇠이다.
\end{formula}

\begin{theorem}[Constant-Coefficient Second-Order ODE]
\[
ay''+by'+cy=0
\]
에는 특성방정식
\[
ar^2+br+c=0
\]
을 쓴다. 근이 $r_1\ne r_2$이면 $C_1e^{r_1x}+C_2e^{r_2x}$, 중근 $r$이면 $(C_1+C_2x)e^{rx}$, $\alpha\pm i\beta$이면
\[
e^{\alpha x}(C_1\cos\beta x+C_2\sin\beta x).
\]
\end{theorem}

\begin{formula}[Separable and First-Order Linear ODEs]
\[
y'=g(x)h(y)\quad\Longrightarrow\quad \int\frac{\dd y}{h(y)}=\int g(x)\,\dd x,
\]
\[
y'+p(x)y=q(x),\qquad
\mu(x)=e^{\int p(x)\dd x},\qquad
(\mu y)'=\mu q.
\]
\end{formula}

\begin{formula}[Constant-Coefficient Systems]
\[
\mathbf x'(t)=\mathbf A\mathbf x(t)
\quad\Longrightarrow\quad
\mathbf x(t)=e^{t\mathbf A}\mathbf x(0).
\]
대각화 가능하면 $e^{\lambda t}\mathbf v$ 꼴의 고유모드들의 선형결합으로 본다.
\end{formula}


\begin{fact}[Three Fundamental PDEs]
\[
\text{Heat: }u_t=\kappa\Delta u,
\qquad
\text{Wave: }u_{tt}=c^2\Delta u,
\qquad
\text{Laplace: }\Delta u=0.
\]
열방정식은 확산, 파동방정식은 파동 전파, Laplace 방정식은 조화함수와 평형을 대표한다.
\end{fact}


\finalchapter{5}{Combinatorics and Discrete Mathematics}

\begin{formula}[Basic Counting]
\[
{}_nP_r=\frac{n!}{(n-r)!},
\qquad
\binom nr=\frac{n!}{r!(n-r)!},
\qquad
\left(\!\binom nr\!\right)=\binom{n+r-1}{r}.
\]
기본식은 이 정도만 남기고, 이후에는 포함배제ㆍCatalanㆍBurnside 같은 구조를 우선 본다.
\end{formula}

\begin{theorem}[Binomial Theorem]
\[
(x+y)^n=\sum_{k=0}^{n}\binom nkx^{n-k}y^k.
\]
특히
\[
\sum_{k=0}^{n}\binom nk=2^n,
\qquad
\sum_{k=0}^{n}(-1)^k\binom nk=0\quad(n\ge1).
\]
\end{theorem}

\begin{formula}[Inclusion--Exclusion]
두 집합:
\[
|A\cup B|=|A|+|B|-|A\cap B|.
\]
세 집합:
\[
|A\cup B\cup C|
=\sum|A|-\sum|A\cap B|+|A\cap B\cap C|.
\]
``적어도 하나''를 직접 세기 복잡하면 여사건를 먼저 본다.
\end{formula}

\begin{theorem}[Pigeonhole Principle]
$n+1$개의 대상을 $n$개의 상자에 넣으면 어떤 상자에는 최소 두 개가 들어간다. 일반형: $N$개를 $k$개에 넣으면 어떤 상자에는 최소
\[
\left\lceil\frac Nk\right\rceil
\]
개가 들어간다.
\end{theorem}

\begin{formula}[Lattice Paths]
$(0,0)$에서 $(m,n)$까지 right/up만 허용하면 경로 수는
\[
\binom{m+n}{m}.
\]
대각선을 넘지 않는 조건이 붙으면 reflection principle 또는 Catalan number를 의심한다.
\end{formula}

\begin{formula}[Derangements]
$D_n$을 아무 원소도 제자리에 있지 않는 permutations의 수라 하면
\[
D_n=n!\sum_{k=0}^{n}\frac{(-1)^k}{k!},
\qquad
D_n=\left\lfloor\frac{n!}{e}+\frac12\right\rfloor.
\]
\end{formula}

\begin{formula}[Vandermonde and Stars and Bars]
\[
\sum_k\binom{r}{k}\binom{s}{n-k}=\binom{r+s}{n}.
\]
음이 아닌 정수해를 갖는
\[
x_1+\cdots+x_m=n
\]
의 수는 $\binom{n+m-1}{m-1}$이고, 양의 정수해의 수는 $\binom{n-1}{m-1}$이다.
\end{formula}


\begin{formula}[Catalan Numbers]
\[
\boxed{C_n=\frac1{n+1}\binom{2n}{n}}.
\]
Dyck 경로, 올바른 괄호열, 이진트리 계수 등 여러 문제에 반복해서 나타난다.
\end{formula}

\begin{theorem}[Hall's Marriage Theorem]
이분그래프에서 왼쪽 모든 정점을 서로 다른 오른쪽 정점과 매칭할 수 있을 필요충분조건은 모든 부분집합 $S$에 대해
\[
\boxed{|N(S)|\ge|S|}
\]
이다.
\end{theorem}

\begin{fact}[Ramsey Anchor]
가장 유명한 정확한 Ramsey number:
\[
\boxed{R(3,3)=6}.
\]
6명이 있으면 서로 모두 아는 3명 또는 서로 모두 모르는 3명이 반드시 존재한다는 고전적 해석이 있다.
\end{fact}

\begin{theorem}[Burnside's Lemma]
유한군 $G$가 유한집합 $X$에 작용할 때 orbit의 개수는
\[
|X/G|=\frac1{|G|}\sum_{g\in G}|\operatorname{Fix}(g)|.
\]
``회전/대칭을 같은 것으로 본다''는 counting 문제에서 떠올린다.
\end{theorem}

\begin{formula}[Linear Recurrence]
$a_n=c_1a_{n-1}+\cdots+c_ka_{n-k}$에는 특성다항식을 쓴다. Fibonacci의 특성방정식은
\[
r^2-r-1=0.
\]
\end{formula}

\begin{theorem}[Master Theorem for Divide-and-Conquer]
\[
T(n)=aT(n/b)+f(n),
\qquad d=\log_ba.
\]
표준 세 case만 기억한다.
\begin{enumerate}
  \item $f(n)=O(n^{d-\varepsilon})$: $T(n)=\Theta(n^d)$.
  \item $f(n)=\Theta(n^d(\log n)^k)$: $T(n)=\Theta(n^d(\log n)^{k+1})$.
  \item $f(n)=\Omega(n^{d+\varepsilon})$이고 정칙성 조건이 성립: $T(n)=\Theta(f(n))$.
\end{enumerate}
서로 다른 크기의 부분문제에는 그대로 적용하지 않는다.
\end{theorem}

\begin{theorem}[Handshaking Lemma]
그래프에서
\[
\sum_{v\in V}\deg(v)=2|E|.
\]
따라서 홀수 차수 꼭짓점의 개수는 항상 짝수이다.
\end{theorem}

\begin{theorem}[Eulerian Trail Criterion]
연결 그래프에서 모든 꼭짓점 차수가 짝수이면 Euler 회로가 존재한다. 정확히 두 꼭짓점만 홀수 차수이면 그 두 점을 끝점으로 하는 Euler 경로가 존재한다. K\"onigsberg 다리 문제의 핵심이다.
\end{theorem}

\begin{formula}[Trees]
$n$ 꼭짓점을 가진 tree는 정확히
\[
|E|=n-1
\]
개의 변을 가진다. 연결 그래프에서 cycle이 없다는 것과 $|E|=|V|-1$은 tree의 대표적 동치 조건이다.
\end{formula}

\begin{theorem}[Planar Graph Euler Formula]
연결 평면그래프의 평면 embedding에서
\[
V-E+F=2.
\]
단순 평면그래프 $(V\ge3)$이면
\[
E\le3V-6,
\]
삼각형을 포함하지 않으면 $E\le2V-4$.
\end{theorem}

\begin{fact}[Spanning Trees]
트리는 연결되고 순환이 없는 그래프이며, $n$개의 꼭짓점을 가지면 변은 정확히 $n-1$개이다. 완전그래프의 신장트리 수는 Cayley 공식
\[
\boxed{T(K_n)=n^{n-2}}.
\]
\end{fact}

\begin{formula}[Moser's Circle Problem]
원 위 $n$점을 모든 현으로 잇고 내부에서 세 현이 한 점에 만나지 않으면 최대 영역 수는
\[
\boxed{1+\binom n2+\binom n4}.
\]
첫 값은 $1,2,4,8,16,31,\ldots$로 $n=6$에서 $2^{n-1}$ 패턴이 깨진다.
\end{formula}

\begin{fact}[Eulerian vs Hamiltonian]
Euler 경로는 \textbf{모든 변}을 정확히 한 번, Hamilton 경로는 \textbf{모든 꼭짓점}을 정확히 한 번 방문한다. 이름을 바꿔 쓰지 않는다.
\end{fact}

\begin{formula}[Tower of Hanoi]
$n$개의 원판을 옮기는 최소 이동 횟수는
\[
\boxed{2^n-1}.
\]
점화식 $T_n=2T_{n-1}+1$의 대표 anchor이다.
\end{formula}


\begin{theorem}[Euler's Polyhedron Formula]
Convex polyhedron에 대해
\[
V-E+F=2.
\]
각 꼭짓점의 차수가 $3$이고 $f_k$가 $k$각형 면의 수이면
\[
\sum_{k\ge3}(6-k)f_k=12.
\]
따라서 모든 face가 pentagon 또는 hexagon이면 pentagons는 정확히 $12$개이다.
\end{theorem}

\finalheading{경우의 수 문제를 받으면}
순서가 중요한가, 중복이 허용되는가, 여사건이 쉬운가, 포함배제이 필요한가, 격자경로로 바뀌는가, 대각선 제약가 Catalan/reflection으로 가는가를 순서대로 본다.


\finalchapter{6}{Probability and Statistics}

\begin{formula}[Conditional Probability and Bayes]
\[
\Prob(A\mid B)=\frac{\Prob(A\cap B)}{\Prob(B)}.
\]
Bayes' theorem:
\[
\Prob(A_i\mid B)=
\frac{\Prob(B\mid A_i)\Prob(A_i)}{\sum_j\Prob(B\mid A_j)\Prob(A_j)}.
\]
조건부확률에서는 sample space가 바뀐다는 점을 먼저 생각한다.
\end{formula}

\begin{definition}[Independence]
\[
A\perp B
\quad\Longleftrightarrow\quad
\Prob(A\cap B)=\Prob(A)\Prob(B).
\]
서로소(disjoint)와 독립(independent)은 다르다. 양의 확률을 갖는 서로소 사건은 독립일 수 없다.
\end{definition}

\begin{formula}[Expectation and Variance]
이산확률변수:
\[
\Exp[X]=\sum_xx\Prob(X=x).
\]
\[
\Var(X)=\Exp[X^2]-\Exp[X]^2.
\]
독립인 $X,Y$에 대해
\[
\Var(X+Y)=\Var(X)+\Var(Y).
\]
\end{formula}

\begin{theorem}[Linearity of Expectation]
\[
\Exp\left[\sum_{i=1}^n a_iX_i\right]
=\sum_{i=1}^n a_i\Exp[X_i].
\]
\boxed{\text{Independence는 필요 없다.}}
개수의 기댓값을 묻는 문제에서는 indicator random variables를 먼저 고려한다.
\end{theorem}

\begin{proposition}[Product of Independent Variables]
독립인 적분가능 확률변수에 대해
\[
\Exp[X_1\cdots X_n]=\prod_{i=1}^n\Exp[X_i].
\]
합의 기댓값과 달리 곱의 기댓값에서는 독립성이 중요하다.
\end{proposition}

\begin{formula}[Total Probability and Total Expectation]
분할 $B_1,\dots,B_m$에 대해
\[
\Prob(A)=\sum_i\Prob(A\mid B_i)\Prob(B_i),
\]
\[
\Exp[X]=\Exp\!\left[\Exp[X\mid Y]\right].
\]
또한
\[
\Var(X)=\Exp[\Var(X\mid Y)]+\Var(\Exp[X\mid Y]).
\]
\end{formula}

\begin{formula}[Covariance and Independence]
\[
\Cov(X,Y)=\Exp[XY]-\Exp[X]\Exp[Y],
\]
\[
\Var(X+Y)=\Var X+\Var Y+2\Cov(X,Y).
\]
독립이면 covariance는 $0$이지만 converse는 일반적으로 거짓이다.
\end{formula}

\begin{fact}[Memoryless Distributions]
Geometric과 exponential distribution은 무기억성(memoryless)을 가진다:
\[
\Prob(X>s+t\mid X>s)=\Prob(X>t).
\]
이산형과 연속형의 대표적인 memoryless 분포이다.
\end{fact}


\begin{formula}[Tail-Sum Formula]
$X\in\Nzero$이면
\[
\boxed{\Exp[X]=\sum_{k=1}^{\infty}\Prob(X\ge k)}.
\]
Waiting-time 문제에서는 ``아직 끝나지 않았을 확률''을 더하는 방식으로 기억한다.
\end{formula}

\begin{formula}[Common Distributions: pmf/pdf, Mean, Variance]
\begin{center}
\small
\renewcommand{\arraystretch}{1.15}
\begin{tabularx}{0.98\textwidth}{@{}L{0.20\textwidth}Y L{0.16\textwidth}L{0.16\textwidth}@{}}
\toprule
Distribution & pmf / pdf & $\Exp[X]$ & $\Var(X)$\\
\midrule
$\Bern(p)$ & $\Prob(X=x)=p^x(1-p)^{1-x}$, $x=0,1$ & $p$ & $p(1-p)$\\
$\Bin(n,p)$ & $\binom nkp^k(1-p)^{n-k}$ & $np$ & $np(1-p)$\\
$\Geom(p)$ & $(1-p)^{k-1}p$, $k\in\N$ & $1/p$ & $(1-p)/p^2$\\
$\Poi(\lambda)$ & $e^{-\lambda}\lambda^k/k!$ & $\lambda$ & $\lambda$\\
$\Unif(a,b)$ & $1/(b-a)$ on $[a,b]$ & $(a+b)/2$ & $(b-a)^2/12$\\
$\ExpDist(\lambda)$ & $\lambda e^{-\lambda x}$, $x\ge0$ & $1/\lambda$ & $1/\lambda^2$\\
$\Normal(\mu,\sigma^2)$ & $\frac1{\sigma\sqrt{2\pi}}e^{-(x-\mu)^2/(2\sigma^2)}$ & $\mu$ & $\sigma^2$\\
\bottomrule
\end{tabularx}
\end{center}
Geometric distribution은 MATH000 convention대로 \textbf{첫 성공이 나온 시행까지 센다}.
\end{formula}




\begin{formula}[Markov and Chebyshev Inequalities]
$X\ge0$이면
\[
\boxed{\Prob(X\ge a)\le\frac{\Exp[X]}a}.
\]
평균 $\mu$, 분산 $\sigma^2$이면
\[
\boxed{\Prob(|X-\mu|\ge t)\le\frac{\sigma^2}{t^2}}.
\]
\end{formula}

\begin{theorem}[Law of Large Numbers and Central Limit Theorem]
독립 동일분포, 평균 $\mu$, 분산 $\sigma^2<\infty$인 $X_i$에 대해
\[
\bar X_n\to\mu\quad\text{(확률적으로; LLN)},
\]
\[
\boxed{\frac{\sqrt n(\bar X_n-\mu)}{\sigma}\Rightarrow\Normal(0,1)}\quad\text{(CLT)}.
\]
정확한 조건보다 ``평균은 $\mu$로 모이고, 중심화ㆍ정규화하면 정규분포''라는 결론을 우선 기억한다.
\end{theorem}

\begin{formula}[Stationary Distribution]
Markov chain의 전이행렬이 $\mathbf P$일 때 stationary distribution $\boldsymbol\pi$는
\[
\boxed{\boldsymbol\pi\mathbf P=\boldsymbol\pi},
\qquad
\sum_i\pi_i=1.
\]
2-state 문제는 이 연립방정식으로 즉시 푼다.
\end{formula}

\begin{formula}[Maxima and Minima]
독립 동일분포인 $X_1,\ldots,X_n$의 CDF가 $F$이면
\[
\Prob\!\left(\max_iX_i\le x\right)=F(x)^n,
\]
\[
\Prob\!\left(\min_iX_i>x\right)=(1-F(x))^n.
\]
Max/min은 CDF로 바꾸면 곱셈으로 정리된다.
\end{formula}

\begin{formula}[Birthday Problem]
$n$명이 $365$일에 균등하게 태어난다고 단순화하면 모두 다른 생일일 확률은
\[
\prod_{k=0}^{n-1}\frac{365-k}{365},
\]
따라서 생일이 같은 pair가 적어도 하나 있을 확률은
\[
1-\prod_{k=0}^{n-1}\frac{365-k}{365}.
\]
$n=23$에서 이미 $1/2$를 넘는다.
\end{formula}

\begin{fact}[Three Famous Probability Anchors]
\begin{itemize}
  \item \textbf{Coupon collector}: $n$종류를 모두 모을 때까지의 기댓값은 $\boxed{nH_n\sim n\log n}$.
  \item \textbf{Broken stick}: 선분을 독립 균등한 두 점에서 자른 세 조각이 삼각형을 이룰 확률은 $\boxed{1/4}$.
  \item \textbf{Buffon's needle}: 간격 $d$, 길이 $\ell\le d$이면 선을 가로지를 확률은 $\boxed{2\ell/(\pi d)}$.
\end{itemize}
\end{fact}

\begin{fact}[Monty Hall]
문 하나를 고른 뒤 진행자가 항상 goat가 있는 다른 문 하나를 열고 switch를 허용하는 표준 Monty Hall 문제에서
\[
\boxed{\Prob(\text{stay wins})=\frac13,
\qquad
\Prob(\text{switch wins})=\frac23.}
\]
\end{fact}

\begin{formula}[Expected Number of Runs]
독립 Bernoulli 수열에서 run의 개수는 인접한 값의 변화 지시변수로 쓴다. 공정한 동전을 $n$번 던지면
\[
\Exp[R_n]=1+\sum_{i=2}^{n}\Prob(X_i\ne X_{i-1})
=1+\frac{n-1}{2}=\frac{n+1}{2}.
\]
\end{formula}

\begin{fact}[Two Stopping-Time Anchors from 2025]
\begin{itemize}
  \item 공정한 동전을 두 번 연속 같은 결과가 나올 때까지 던지면 기대 시행횟수는 $\boxed{3}$.
  \item 공정한 주사위를 $6$이 나올 때까지 던져 나온 눈의 합을 더하면 기대합은 $\boxed{21}$.
\end{itemize}
두 문제 모두 ``상태를 두고 기댓값 점화식을 세운다''는 공통 핵심을 가진다.
\end{fact}

\finalheading{기댓값 문제를 받으면}
\begin{enumerate}
  \item ``개수''를 묻는가? $\to$ 지시변수 + 선형성.
  \item 곱인가? $\to$ 독립성 확인.
  \item ``적어도 하나''인가? $\to$ 여사건.
  \item ``~가 나올 때까지 / 처음으로''인가? $\to$ 기하분포, 점화식, 꼬리합 공식.
  \item 조건이 붙었는가? $\to$ 조건부 표본공간을 다시 만든다.
\end{enumerate}


\finalchapter{7}{Geometry}

\finalsection{Synthetic Geometry}

\begin{formula}[Triangle Essentials]
\[
a^2+b^2=c^2\quad\text{(직각삼각형)},
\]
\[
\frac{a}{\sin A}=\frac{b}{\sin B}=\frac{c}{\sin C}=2R,
\]
\[
c^2=a^2+b^2-2ab\cos C.
\]
넓이 공식:
\[
K=\frac12ab\sin C=rs=\frac{abc}{4R},
\]
\[
K=\sqrt{s(s-a)(s-b)(s-c)}.
\]
\end{formula}



\begin{formula}[Circle Essentials]
반지름 $R$인 원에서 중심각 $\theta$ (라디안)에 대해
\[
\text{arc length}=R\theta,\qquad \text{sector area}=\frac12R^2\theta.
\]
접선은 접점에서 반지름과 수직이고, 원 밖의 한 점에서 그은 두 접선의 길이는 같다.
\end{formula}


\begin{theorem}[Thales' Theorem]
지름이 만드는 원주각은 $90^\circ$이다. 원과 직각이 함께 나오면 지름을 먼저 의심한다.
\end{theorem}

\begin{formula}[Ceva and Menelaus]
삼각형 $ABC$에서 $D\in BC$, $E\in CA$, $F\in AB$.
Concurrent 조건:
\[
\frac{BD}{DC}\frac{CE}{EA}\frac{AF}{FB}=1.
\]
일직선 조건(Menelaus)은 유향비 convention 아래 대응하는 곱이 $-1$이다.
\end{formula}

\begin{fact}[Triangle Centers and Euler Line]
무게중심 $G$는 중선의 교점이며 각 중선을 꼭짓점 쪽에서 $2:1$로 나눈다. 외심 $O$, 무게중심 $G$, 수심 $H$는 한 직선 위에 있고
\[
\boxed{OG:GH=1:2}.
\]
내심는 각의 이등분선의 교점이다.
\end{fact}

\begin{fact}[Platonic Solids]
정다면체는 정확히 다섯 개: tetrahedron, cube, octahedron, dodecahedron, icosahedron. 면의 수는 각각
\[
4,\ 6,\ 8,\ 12,\ 20.
\]
\end{fact}

\finalsection{Analytic Geometry}



\begin{formula}[평면 and 점--평면 Distance]
평면
\[
ax+by+cz=d
\]
의 법선벡터는 $(a,b,c)$. 점 $(x_0,y_0,z_0)$에서 평면까지의 거리는
\[
\frac{|ax_0+by_0+cz_0-d|}{\sqrt{a^2+b^2+c^2}}.
\]
구가 평면에 접하면 중심에서 평면까지의 거리가 반지름과 같다.
\end{formula}

\begin{formula}[구]
중심 $\mathbf c$와 반지름 $r$인 sphere:
\[
\|\mathbf x-\mathbf c\|^2=r^2.
\]
좌표평면 $x=0,y=0,z=0$에 모두 접하는 제1옥탄트의 구는 중심이 $(r,r,r)$ 꼴이다.
\end{formula}



\begin{formula}[Shoelace Formula]
꼭짓점 $(x_i,y_i)$가 순서대로 주어진 polygon의 넓이는
\[
\boxed{A=\frac12\left|\sum_i x_i y_{i+1}-y_i x_{i+1}\right|}.
\]
좌표가 주어진 다각형에서는 가장 먼저 확인한다.
\end{formula}

\begin{fact}[Euclidean, Hyperbolic, Spherical]
한 직선 밖의 한 점을 지나는 평행선의 수는 각각
\[
1,\quad \text{무한히 많음},\quad 0
\]
이며, 삼각형 내각합은 각각
\[
180^\circ,\quad <180^\circ,\quad >180^\circ.
\]
Euclidean 경우의 ``정확히 하나''는 Playfair's axiom으로 불린다.
\end{fact}

\begin{fact}[Archimedes' 구 and Cylinder]
반지름 $r$인 구와 이를 꼭 맞게 둘러싼 반지름 $r$, 높이 $2r$인 원기둥에서
\[
V_{\text{sphere}}:V_{\text{cylinder}}=2:3.
\]
Archimedes와 Fields Medal 뒷면 디자인의 연결까지 함께 기억한다.
\end{fact}

\begin{theorem}[Pick's Theorem]
lattice polygon의 넓이는
\[
A=I+\frac B2-1,
\]
여기서 $I$는 내부 격자점, $B$는 경계 격자점의 수이다.
\end{theorem}


\finalchapter{8}{Number Theory}

\begin{formula}[Euclidean Algorithm]
\[
\gcd(a,b)=\gcd(b,a\bmod b).
\]
또한
\[
\boxed{\gcd(a,b)\lcm(a,b)=|ab|}.
\]
\end{formula}

\begin{theorem}[B\'ezout's Identity]
어떤 $x,y\in\Z$가 존재하여
\[
\gcd(a,b)=ax+by.
\]
따라서
\[
ax+by=c\text{가 정수해을 가짐}
\iff
\gcd(a,b)\mid c.
\]
\end{theorem}

\begin{formula}[Prime Factorization and Divisor Functions]
$n=\prod p_i^{a_i}$이면
\[
\boxed{\tau(n)=\prod_i(a_i+1)},
\qquad
\boxed{\sigma(n)=\prod_i\frac{p_i^{a_i+1}-1}{p_i-1}}.
\]
또한
\[
\tau(n)\text{ odd}\iff n\text{ is a perfect square}.
\]
\end{formula}

\begin{formula}[Legendre's Formula]
소수 $p$에 대해
\[
v_p(n!)=\sum_{k\ge1}\left\lfloor\frac{n}{p^k}\right\rfloor.
\]
계승의 나눗셈 가능성과 소인수 지수를 묻는 문제에서 즉시 쓴다.
\end{formula}

\begin{formula}[Congruence]
\[
a\equiv b\pmod n\iff n\mid(a-b).
\]
덧셈ㆍ뺄셈ㆍ곱셈은 그대로 가능하지만 나눗셈은 역원의 존재를 확인한다.
\end{formula}

\begin{theorem}[Linear Congruence]
\[
ax\equiv b\pmod n
\]
이 풀릴 필요충분조건은
\[
\boxed{\gcd(a,n)\mid b}.
\]
$d=\gcd(a,n)$가 $b$를 나누면 modulo $n$에서 정확히 $d$개의 해의 residue classes가 생긴다.
\end{theorem}

\begin{formula}[Modular Inverse]
\[
a^{-1}\pmod n\text{ exists}
\iff
\gcd(a,n)=1.
\]
확장 Euclidean algorithm으로 역원를 실제로 구할 수 있다.
\end{formula}

\begin{fact}[Small-Modulus Square Residues]
\[
x^2\pmod4\in\{0,1\},
\qquad
x^2\pmod8\in\{0,1,4\}.
\]
따라서 두 squares의 합은 modulo $4$에서 $3$이 될 수 없다. 2010 과학퀴즈의 $x^2+y^2=963$ 불가능성의 핵심이다.
\end{fact}

\finalheading{정수해의 존재성을 물으면}
작은 modulus를 먼저 본다. 특히 $2,4,8,3,5,9$에서 square/cube residue가 강하게 제한되면 복잡한 factorization보다 훨씬 빨리 끝난다.

\begin{theorem}[Fermat and Euler]
소수 $p$에서 $p\nmid a$이면
\[
a^{p-1}\equiv1\pmod p.
\]
일반적으로 $\gcd(a,n)=1$이면
\[
a^{\varphi(n)}\equiv1\pmod n.
\]
\end{theorem}

\begin{formula}[Euler's Totient]
\[
\boxed{\varphi(n)=n\prod_{p\mid n}\left(1-\frac1p\right)}.
\]
특히 $p$가 소수이면 $\varphi(p)=p-1$이고
\[
\varphi(p^k)=p^k-p^{k-1}.
\]
\end{formula}

\begin{theorem}[Chinese Remainder Theorem]
$m,n$이 coprime이면
\[
x\equiv a\pmod m,
\qquad
x\equiv b\pmod n
\]
은 modulo $mn$에서 유일한 solution class를 가진다. 서로소인 여러 moduli에도 동일하게 확장된다.
\end{theorem}

\begin{theorem}[Wilson's Theorem]
$p>1$에 대해
\[
p\text{ prime}
\iff
(p-1)!\equiv-1\pmod p.
\]
유명하고 짧은 소수 판정식으로 기억한다.
\end{theorem}




\begin{formula}[Multiplicative Order]
$\gcd(a,n)=1$일 때 $\ord_n(a)$는 $a^k\equiv1\pmod n$인 최소 양의 $k$. 항상
\[
\ord_n(a)\mid\varphi(n).
\]
위수를 알면 큰 지수를 빠르게 줄일 수 있다.
\end{formula}

\begin{fact}[Euler Criterion and Quadratic Residues]
홀수 소수 $p\nmid a$에 대해
\[
a^{(p-1)/2}\equiv\left(\frac ap\right)\pmod p.
\]
이차상호법칙은 홀수 소수 $p,q$에 대해
\[
\left(\frac pq\right)\left(\frac qp\right)=(-1)^{\frac{p-1}{2}\frac{q-1}{2}}.
\]
\end{fact}


\begin{formula}[GCD of Exponential Differences]
$a>1$에 대해
\[
\boxed{\gcd(a^m-1,a^n-1)=a^{\gcd(m,n)}-1}.
\]
지수가 큰 gcd 문제에서 바로 쓴다.
\end{formula}

\begin{theorem}[Prime Number Theorem and Bertrand]
\[
\boxed{\pi(x)\sim\frac{x}{\log x}}.
\]
또한 모든 $n>1$에 대해
\[
\boxed{n<p<2n}
\]
인 소수 $p$가 존재한다 (Bertrand의 공준).
\end{theorem}

\begin{theorem}[Two Squares and Perfect Numbers]
소수 $p$에 대해
\[
p=x^2+y^2\iff p=2\text{ or }p\equiv1\pmod4.
\]
짝수 완전수는 정확히
\[
\boxed{2^{p-1}(2^p-1)}
\]
꼴이며 $2^p-1$이 소수여야 한다 (Euclid--Euler theorem).
\end{theorem}

\begin{fact}[Mih\u{a}ilescu / Catalan]
$A^m-B^n=1$인 $A,B,m,n>1$의 유일한 해는
\[
\boxed{3^2-2^3=1}.
\]
즉 $8$과 $9$가 연속한 유일한 자명하지 않은 완전거듭제곱이다.
\end{fact}

\begin{fact}[Continued Fractions]
단순 연분수는 유리수에서 유한하게 끝난다. 반대로 이차무리수에서는 결국 주기적이다 (Lagrange theorem). 대표적으로
\[
\varphi=[1;1,1,1,\ldots].
\]
\end{fact}

\begin{fact}[Famous Number-Theory Anchors]
\begin{itemize}
  \item Goldbach 추측: 모든 짝수 $>2$는 두 소수의 합인가? 아직 미해결.
  \item 쌍둥이 소수 추측: 차이가 $2$인 소수쌍이 무한히 많은가? 아직 미해결.
  \item 이차상호법칙: 가우스 (Carl Friedrich Gauss)의 대표적인 정수론 결과.
  \item Fermat의 마지막 정리: $n>2$에서 $x^n+y^n=z^n$의 0이 아닌 정수해가 없으며 앤드루 와일스 (Andrew Wiles)가 증명했다.
\end{itemize}
\end{fact}


\finalchapter{9}{Algebra}

\begin{theorem}[Lagrange and Orbit--Stabilizer]
유한군 $G$의 부분군 $H$에 대해
\[
|H|\mid|G|.
\]
군 $G$가 집합 $X$에 작용하면
\[
\boxed{|G\cdot x|=[G:G_x]=\frac{|G|}{|G_x|}}.
\]
따라서 원소의 위수도 $|G|$를 나눈다.
\end{theorem}

\begin{theorem}[First Isomorphism Theorem]
군 준동형 $\varphi:G\to H$에 대해
\[
\boxed{G/\ker\varphi\cong\im\varphi}.
\]
특히 $\varphi$가 단사일 필요충분조건은 $\ker\varphi=\{e\}$이다.
\end{theorem}

\begin{formula}[Cyclic Groups and Homomorphisms]
\[
|\Z/n\Z|=n,
\qquad
\#\{\text{generators of }\Z/n\Z\}=\varphi(n).
\]
가법군 사이에서는
\[
\boxed{|\operatorname{Hom}(\Z/m\Z,\Z/n\Z)|=\gcd(m,n)}.
\]
\end{formula}

\begin{formula}[Permutations]
순열의 위수는 서로소 순환 길이들의 최소공배수이다:
\[
\boxed{\ord(\sigma)=\lcm(\ell_1,\ldots,\ell_r)}.
\]
\[
|S_n|=n!,\qquad |A_n|=\frac{n!}{2}\quad(n\ge2).
\]
\end{formula}

\begin{formula}[Finite Fields and Classical Groups]
유한체는 정확히 소수 거듭제곱 크기에서 존재하며, 주어진 $q=p^r$에 대해 동형까지 유일하다:
\[
\boxed{|\F_q|=q=p^r}.
\]
\[
\boxed{|\GL_n(\F_q)|=\prod_{k=0}^{n-1}(q^n-q^k)}.
\]
$\SL_n(\F_q)$는 행렬식이 $1$인 부분군이다.
\end{formula}

\begin{formula}[Vi\`ete and Root Identities]
\[
x^n+a_{n-1}x^{n-1}+\cdots+a_0=\prod_{i=1}^n(x-r_i)
\]
이면
\[
\sum_i r_i=-a_{n-1},\qquad \prod_i r_i=(-1)^n a_0.
\]
또한 $p(a)\ne0$이면
\[
\boxed{\frac{p'(a)}{p(a)}=\sum_i\frac1{a-r_i}}.
\]
근을 직접 구하지 않고 대칭식을 계산하는 데 강력하다.
\end{formula}

\begin{theorem}[Abel--Ruffini]
5차 이상의 일반 다항방정식에는 사칙연산과 근호만으로 모든 해를 나타내는 보편적 공식이 존재하지 않는다. ``모든 5차방정식이 풀리지 않는다''는 뜻은 아니다.
\end{theorem}

\begin{fact}[Quaternions]
Hamilton의 quaternion $\mathbb H$에서는
\[
i^2=j^2=k^2=ijk=-1,
\qquad ij=k,\quad ji=-k.
\]
대표적인 비가환 나눗셈대수이다.
\end{fact}

\finalchapter{10}{Topology}

\begin{formula}[Basic Notation]
\[
A^\circ=\interior A,\qquad \overline A=\closure A,\qquad \partial A=\overline A\setminus A^\circ.
\]
위상동형사상 (homeomorphism)은 연속인 전단사이고 역함수도 연속인 사상이다. 정의보다 아래의 불변량과 유명 정리를 우선 기억한다.
\end{formula}

\begin{fact}[Compactness and Connectedness]
콤팩트 집합의 연속상은 콤팩트이고, 연결집합의 연속상은 연결이다. 특히
\[
\boxed{K\subset\R^n\text{ compact}\iff K\text{ closed and bounded}.}
\]
연결공간에서 $\Q$처럼 완전비연결인 공간으로 가는 연속함수는 강하게 제한된다.
\end{fact}

\begin{fact}[Fundamental Group Anchors]
\[
\pi_1(S^1)\cong\Z,
\qquad
\pi_1(S^n)=0\quad(n\ge2),
\qquad
\pi_1(T^2)\cong\Z^2.
\]
원의 감김수 (winding number), 구면의 단일연결성, 원환면의 두 독립 루프를 연결해 기억한다.
\end{fact}

\begin{formula}[Euler Characteristic]
\[
\chi=V-E+F,
\qquad
\boxed{\chi(\Sigma_g)=2-2g}
\]
닫힌 가향곡면 $\Sigma_g$에 대해. 따라서 구면은 $2$, 원환면은 $0$이다.
\end{formula}

\begin{theorem}[Three Famous Topological Theorems]
\begin{itemize}
  \item \textbf{Brouwer 고정점정리}: 닫힌 공의 연속 자기함수는 고정점을 갖는다.
  \item \textbf{Hairy Ball 정리}: 짝수차원 구면, 특히 $S^2$에는 어디서도 0이 되지 않는 연속 접벡터장이 없다.
  \item \textbf{Borsuk--Ulam 정리}: 연속함수 $f:S^n\to\R^n$에는 $f(x)=f(-x)$인 대척점 쌍이 존재한다.
\end{itemize}
\end{theorem}

\begin{fact}[M\"obius Strip and Klein Bottle]
M\"obius strip은 비가향이며 경계가 하나 있다. Klein bottle은 닫힌 비가향곡면이고 $\R^3$에 자기교차 없이 매장할 수 없다.
\end{fact}

\begin{fact}[Poincar\'e Conjecture]
닫힌 단일연결 $3$차원 다양체는 $S^3$와 위상동형이다. 그리고리 페렐만 (Grigori Perelman)이 Ricci 흐름을 이용해 해결했고 2006년 필즈상을 거부했다.
\end{fact}

\begin{fact}[Thurston's Eight Geometries]
Thurston의 기하화 이론은 $3$차원 다양체를 여덟 가지 모델 기하와 연결한다. 과학퀴즈에서는
\[
\boxed{\text{Thurston}\longleftrightarrow 8\text{ geometries}\longleftrightarrow 3\text{-manifolds}}
\]
를 우선 기억한다.
\end{fact}

